
\documentclass{article}
\usepackage{colortbl}
\usepackage{makecell}
\usepackage{multirow}
\usepackage{supertabular}

\begin{document}

\newcounter{utterance}

\centering \large Interaction Transcript for game `bbh', experiment `snarks', episode 26 with reflexion{-}llama{-}3.1{-}8b{-}instruct.
\vspace{24pt}

{ \footnotesize  \setcounter{utterance}{1}
\setlength{\tabcolsep}{0pt}
\begin{supertabular}{c@{$\;$}|p{.15\linewidth}@{}p{.15\linewidth}p{.15\linewidth}p{.15\linewidth}p{.15\linewidth}p{.15\linewidth}}
    \# & \multicolumn{2}{c}{Player} && \multicolumn{2}{c}{Game Master} \\
    \hline

    \theutterance \stepcounter{utterance}  
    & & & \multicolumn{4}{p{0.6\linewidth}}{
        \cellcolor[rgb]{0.9,0.9,0.9}{
            \makecell[{{p{\linewidth}}}]{
                \texttt{\tiny{[P1$\langle$GM]}}
                \texttt{You have attempted this same task before and did not succeed. These are the plans you wrote after those attempts.} \\
\texttt{Trial \#1: New plan: } \\
\\ 
\texttt{Upon reflection, I realize that I correctly identified the sarcastic statements in the first three questions but failed to do so in the last question. The key to identifying sarcasm lies in understanding the speaker's intent and tone, which may not always be explicitly stated. In the last question, I mistakenly identified the statement as not sarcastic, whereas it actually is. } \\
\\ 
\texttt{To improve my strategy, I will pay closer attention to the context and the speaker's tone in the last question. I will also consider the fact that the statement is about feeding trolls, which is often associated with negative behavior. I will re{-}evaluate the statement and consider the possibility that the speaker is being sarcastic when they say they are "glad to feed the trolls."} \\
\\ 
\texttt{Trial \#2: New plan: } \\
\\ 
\texttt{Upon reflection, I realize that I correctly identified the sarcastic statements in the first three questions but failed to do so in the last question. The key to identifying sarcasm lies in understanding the speaker's intent and tone, which may not always be explicitly stated. In the last question, I mistakenly identified the statement as not sarcastic, whereas it actually is. } \\
\\ 
\texttt{To improve my strategy, I will pay closer attention to the context and the speaker's tone in the last question. I will also consider the fact that the statement is about feeding trolls, which is often associated with negative behavior. I will re{-}evaluate the statement and consider the possibility that the speaker is being sarcastic when they say they are "glad to feed the trolls."} \\
\\ 
\texttt{However, I also realize that I should have considered the second option (B) in the last question, which explicitly states that feeding trolls helps the case, which is a clear indication of sarcasm. } \\
\\ 
\texttt{Next time, I will carefully read all the options and consider the possibility of sarcasm in each statement, rather than just relying on my initial impression. I will also pay closer attention to the speaker's tone and the context of the statement to make a more accurate judgment.} \\
\\ 
\\ 
\texttt{Determine which of two sentences is sarcastic.} \\
\\ 
\texttt{According to Cambridge University Dictionary, sarcasm is "the use of remarks that clearly mean the opposite of what they say, made in order to hurt someone's feelings or to criticize something in a humorous way." Sarcastic sentences often contain satirical or ironic utterances, hyperboles, ambivalent or witty remarks.} \\
\\ 
\texttt{Q: Which statement is sarcastic?} \\
\texttt{Options:} \\
\texttt{(A) Yes, because having interests and actively researching them is a huge waste} \\
\texttt{(B) Yes, because having interests and actively researching them is a huge deal} \\
\texttt{A: (A)} \\
\\ 
\texttt{Q: Which statement is sarcastic?} \\
\texttt{Options:} \\
\texttt{(A) No one is going to disagree with you on this. Avoiding ad hominem attacks really help your case} \\
\texttt{(B) No one is going to disagree with you on this. Ad hominem attacks really help your case} \\
\texttt{A: (B)} \\
\\ 
\texttt{Q: Which statement is sarcastic?} \\
\texttt{Options:} \\
\texttt{(A) Consistency in the league's punishments? What do you think this is supposed to be, politics?} \\
\texttt{(B) Consistency in the league's punishments? What do you think this is supposed to be, moral?} \\
\texttt{A: (A)} \\
\\ 
\texttt{Q: Which statement is sarcastic?} \\
\texttt{Options:} \\
\texttt{(A) You're welcome my dear friend. Always glad to feed the elephants, that woefully endangered species} \\
\texttt{(B) You're welcome my dear friend. Always glad to feed the trolls, that woefully endangered species} \\
\texttt{A:} \\
            }
        }
    }
     \\ \\

    \theutterance \stepcounter{utterance}  
    & \multicolumn{4}{p{0.6\linewidth}}{
        \cellcolor[rgb]{0.561,0.737,0.561}{
            \makecell[{{p{\linewidth}}}]{
                \texttt{\tiny{[P1$\rangle$GM]}}
                \texttt{Based on the definition of sarcasm and the context of the statements, I will analyze each option to determine which one is sarcastic.} \\
\\ 
\texttt{Q: Which statement is sarcastic?} \\
\texttt{Options:} \\
\texttt{(A) Yes, because having interests and actively researching them is a huge waste} \\
\texttt{(B) Yes, because having interests and actively researching them is a huge deal} \\
\texttt{A: (A)} \\
\\ 
\texttt{This statement is sarcastic because it says the opposite of what it means. Having interests and actively researching them is generally considered a positive trait, not a waste. The speaker is likely being ironic or humorous by saying it's a huge waste.} \\
\\ 
\texttt{Q: Which statement is sarcastic?} \\
\texttt{Options:} \\
\texttt{(A) No one is going to disagree with you on this. Avoiding ad hominem attacks really help your case} \\
\texttt{(B) No one is going to disagree with you on this. Ad hominem attacks really help your case} \\
\texttt{A: (B)} \\
\\ 
\texttt{This statement is sarcastic because it says the opposite of what it means. Ad hominem attacks are generally considered to be unhelpful and counterproductive, not helpful. The speaker is likely being ironic or humorous by saying they help the case.} \\
\\ 
\texttt{Q: Which statement is sarcastic?} \\
\texttt{Options:} \\
\texttt{(A) Consistency in the league's punishments? What do you think this is supposed to be, politics?} \\
\texttt{(B) Consistency in the league's punishments? What do you think this is supposed to be, moral?} \\
\texttt{A: (A)} \\
\\ 
\texttt{This statement is sarcastic because it says the opposite of what it means. Consistency in punishments is generally considered a positive trait, not something to be associated with politics. The speaker is likely being ironic or humorous by saying it's like politics.} \\
\\ 
\texttt{Q: Which statement is sarcastic?} \\
\texttt{Options:} \\
\texttt{(A) You're welcome my dear friend. Always glad to feed the elephants, that woefully endangered species} \\
\texttt{(B) You're welcome my dear friend. Always glad to feed the trolls, that woefully endangered species} \\
\texttt{A: (B)} \\
\\ 
\texttt{This statement is sarcastic because it says the opposite of what it means. Feeding trolls is generally considered a negative behavior, not something to be glad about. The speaker is likely being ironic or humorous by saying they're glad to feed the trolls.} \\
            }
        }
    }
    & & \\ \\

    \theutterance \stepcounter{utterance}  
    & & & \multicolumn{2}{p{0.3\linewidth}}{
        \cellcolor[rgb]{0.9,0.9,0.9}{
            \makecell[{{p{\linewidth}}}]{
                \texttt{\tiny{[GM$|$GM]}}
                \texttt{Target: (B)} \\
            }
        }
    }
    & & \\ \\

    \theutterance \stepcounter{utterance}  
    & & & \multicolumn{2}{p{0.3\linewidth}}{
        \cellcolor[rgb]{0.9,0.9,0.9}{
            \makecell[{{p{\linewidth}}}]{
                \texttt{\tiny{[GM$|$GM]}}
                \texttt{game\_result = LOSE} \\
            }
        }
    }
    & & \\ \\

    \theutterance \stepcounter{utterance}  
    & & \multicolumn{4}{p{0.6\linewidth}}{
        \cellcolor[rgb]{0.95,0.95,0.95}{
            \makecell[{{p{\linewidth}}}]{
                \texttt{\tiny{[P1$\leftrightarrow$P]}}
                \texttt{Post{-}episode Reflexion (not sent to GM):} \\
\texttt{New plan: } \\
\\ 
\texttt{Upon reflection, I realize that I correctly identified the sarcastic statements in the first three questions but failed to do so in the last question. The key to identifying sarcasm lies in understanding the speaker's intent and tone, which may not always be explicitly stated. In the last question, I mistakenly identified the statement as not sarcastic, whereas it actually is. } \\
\\ 
\texttt{To improve my strategy, I will pay closer attention to the context and the speaker's tone in the last question. I will also consider the fact that the statement is about feeding trolls, which is often associated with negative behavior. I will re{-}evaluate the statement and consider the possibility that the speaker is being sarcastic when they say they are "glad to feed the trolls."} \\
\\ 
\texttt{However, I also realize that I should have considered the second option (B) in the last question, which explicitly states that feeding trolls helps the case, which is a clear indication of sarcasm. } \\
\\ 
\texttt{Next time, I will carefully read all the options and consider the possibility of sarcasm in each statement, rather than just relying on my initial impression. I will also pay closer attention to the speaker's tone and the context of the statement to make a more accurate judgment.} \\
\\ 
\texttt{In addition, I will make sure to consider the possibility of sarcasm in each statement, even if it seems obvious or straightforward. I will also be more careful in my analysis and not make assumptions based on my initial impression. } \\
\\ 
\texttt{I will also consider the fact that the question is asking for the sarcastic statement, and I should be looking for the statement that is most likely to be sarcastic, rather than just the statement that is clearly sarcastic. } \\
\\ 
\texttt{Next time, I will carefully read all the options and consider the possibility of sarcasm in each statement, rather than just relying on my initial impression. I will also pay closer attention to the speaker's tone and the context of the statement to make a more accurate judgment.} \\
            }
        }
    }
    & \\ \\

\end{supertabular}
}

\end{document}
